\subsection{Exact transport from the actual character to the standard character}
\label{sec:archimedean-character-transport}

The canonical root-Haar scale \(\theta=1\) does not make the
actual character parameter equal to one. After the prescribed
phase rotation, the actual local datum has
\[
       c=-2ia,\qquad
       a=\mathrm{kaRe}(D)
         =\frac{D.\mathrm{host.ka}}{g_0g_1}>0 .
\]
The parameter \(a\) and the root-Haar scale \(\theta\) are
different quantities. We now transport this datum to \(c_1=-2i\)
and calculate the pointwise and radial factors together.
Every factor whose exponent contains \(k\) cancels exactly.
The calculation uses the original root, quotient and radial
normalizations.\cite{current-construction-record}

\begin{proposition}[Character transport without a compact-rank loss]
\label{prop:archimedean-character-transport}
Let \(s=r+1\), \(k\ge1\), and let \(L\) be an actual local
Bessel datum with character \(c=-2ia\), \(a>0\), and the
period, smooth-decay and all-cell factorization hypotheses
already established. Let \(F\) be a finite compact-finite
coefficient of its discrete series. There is a standard-character
datum \(L_1\), with the same discrete series and compact
parameters, and a single finite compact-finite coefficient
\(F_1\), chosen before all points and radial sections, with
the following properties. The pointwise identities below apply
to every relevant original point on the nonzero-Gram locus
\(\mathcal G_L(\gamma)>0\); the radial identity uses the
original big cell up to its null boundary.

Put
\[
 \eta=a^{-1},\qquad \tau=a^r,\qquad
 t_j=\tau\eta^{r-j}=a^j\quad(0\le j\le r).
                                                        \tag{CT1}
\]
The new negative vectors are \(z_1=\tau z\), \(z'_1=\tau z'\);
thus \(T_1=\tau^2T\), \(T'_1=\tau^2T'\).
The prescribed point transform \(\gamma\mapsto\gamma_1\)
has Gram matrix
\[
       Z_1=\operatorname{diag}(t_j)\,
                       Z\,\operatorname{diag}(t_j).
                                                        \tag{CT2}
\]
Let \(P_\gamma\) be the original root pattern, and write
\[
 d_\gamma=|P_\gamma|,\qquad
 S_\gamma=\sum_{(i,j)\in P_\gamma}(j-i),\qquad
 \alpha_\gamma=\eta^{d_\gamma-2S_\gamma}.
\]
Define the positive numbers
\[
 J=\tau^{4s}\eta^{2rs}=a^{2rs},\qquad
 W=\prod_{j=0}^r
        t_j^{\,4(2(r-j)+k-1)} .
                                                        \tag{CT3}
\]
With the original notation for point weight, orbital square,
Gram polynomial and squared radial norm, one has
\begin{align}
 \mathrm{orbSq}_{L_1}(F_1,\gamma_1)
      &=\alpha_\gamma^{-2}
                         \mathrm{orbSq}_L(F,\gamma),
                                                        \label{eq:character-orb}\\
 \mathrm{ptWeight}_{L_1}(\gamma_1)
      &=W\alpha_\gamma^2\eta^{-2d_\gamma}
                             \mathrm{ptWeight}_L(\gamma),
                                                        \label{eq:character-weight}\\
 \mathcal N_{L_1}(F_1)&=JW\mathcal N_L(F),
                                                        \label{eq:character-norm}\\
 \mathcal G_{L_1}(\gamma_1)&=J\mathcal G_L(\gamma).
                                                        \label{eq:character-gram}
\end{align}
These identities include every admitted quotient measure
and radial section, transported as described in the proof.
In particular a standard-character estimate
\[
 \mathrm{ptWeight}_{L_1}(\gamma_1)
       \mathrm{orbSq}_{L_1}(F_1,\gamma_1)
 \le B(\mathcal G_{L_1}(\gamma_1))\mathcal N_{L_1}(F_1)
\]
implies the actual-character estimate
\[
 \mathrm{ptWeight}_L(\gamma)\mathrm{orbSq}_L(F,\gamma)
 \le a^{2rs-2d_\gamma}
               B(a^{2rs}\mathcal G_L(\gamma))\mathcal N_L(F).
                                                        \tag{CT4}
\]
No division by a radial norm or an orbital value is required.
\end{proposition}

\begin{proof}
Let \(D\) act on the first isotropic basis vector of index
\(j<r\) by \(t_j\), on the opposite vector by \(t_j^{-1}\),
and on the middle space by the identity. This is an element
of the original unitary group. Put
\[
                  \chi(g)=D^{-1}gD .
                                                        \tag{CT5}
\]
The last frame vector is scaled separately, by \(\tau\).
Then
\[
 Df_j(z_1)=t_jf_j(z)
\]
also at \(j=r\). Taking Hermitian pairings proves (CT2).
The elementary root-coordinate calculation gives
\[
 \ell_{z_1}(\chi(n))=\eta^{-1}\ell_z(n).
                                                        \tag{CT6}
\]
Consequently the new character \(c_1=\eta c=-2i\) satisfies
\(\psi_1(\chi(n))=\psi_0(n)\).
Conjugation normalizes the unipotent subgroup and preserves
its Haar class. Give the new unipotent and compact measures
their exact pushforwards under \(\chi\); their positive
scalar normalizations are retained. The group Haar measure
is unchanged by this inner automorphism. The compact
parameters and the discrete series are unchanged.

Initially set \(\widetilde F=F\circ\chi^{-1}\).
This is a finite smooth coefficient of the same representation:
if \(F(g)=P\sigma(g)Q\), then
\[
 \widetilde F(g)
       =(P\sigma(D))\,\sigma(g)\,(\sigma(D^{-1})Q).
                                                        \tag{CT7}
\]
Smoothness is preserved by the representation action.
Conjugation need not preserve finiteness under the chosen
maximal compact subgroup, so that assertion is not used.
To supply the literal compact-finite consumer, use the
period-normalized maps. Exact substitution in their period
identities gives
\[
          B_1=B_0\sigma(D),\qquad B'_1=B'_0\sigma(D).
                                                        \tag{CT8}
\]
If either period map is zero, the all-cell factorization makes
the orbital evaluations and radial norm zero; take \(F_1=0\).
Otherwise these maps are nonzero, and
\(s_{\sigma,1}(\widetilde F)=s_{\sigma,0}(F)\).
Choose one compact-finite rank-one coefficient \(H\)
with \(s_{\sigma,1}(H)\ne0\), as in the proof of
Theorem~\ref{thm:flag-exact-norm-comparison}, and put
\[
 F_1=\frac{s_{\sigma,1}(\widetilde F)}
                         {s_{\sigma,1}(H)}H .
                                                        \tag{CT9}
\]
The all-cell factorization
Theorem~\ref{thm:all-cell-relative-factorization}
shows that \(F_1\) and \(\widetilde F\) have identical
orbital evaluations at every relevant nonzero-Gram point
and identical radial norms. If the numerator in (CT9)
is zero, take \(F_1=0\); the same factorization proves
all these evaluations vanish. This choice is independent
of the point and section.

We next track the canonical root measure, still at
root-Haar scale \(1\). The root pattern is unchanged by
(CT2). In its complex coordinate indexed by \((i,j)\),
conjugation scales the coordinate by \(t_j/t_i\).
The prescribed density per complex coordinate is
proportional to the absolute value of the character
parameter. Thus the new canonical root Haar measure is
\[
 \mu_{R,1}
      =\alpha_\gamma\,\chi_*\mu_{R,0},\qquad
 \alpha_\gamma
      =\eta^{d_\gamma}
           \prod_{(i,j)\in P_\gamma}(t_j/t_i)^2
      =\eta^{d_\gamma-2S_\gamma}.
                                                        \tag{CT10}
\]
This scalar is forced by the root-coordinate Jacobian;
it is not a choice of a new free normalization.
The defining Weil identity therefore gives the new
orbital quotient measure as
\(\alpha_\gamma^{-1}\) times the pushforward of the old one.
Substitution in the absolutely convergent orbital integral,
using (CT6) and (CT7), gives
\[
       \mathrm{OrbInt}_{L_1}(\widetilde F,\chi(g))
             =\alpha_\gamma^{-1}\mathrm{OrbInt}_L(F,g).
                                                        \tag{CT11}
\]
The transform and its inverse give a bijection on all
admitted representatives, root measures and orbital quotient
measures. Taking squared Hilbert--Schmidt norms and then
the defining supremum proves \eqref{eq:character-orb};
(CT9) allows replacement by \(F_1\).

For clarity, the point-density identity can be checked
block by block. A block of length \(h\), beginning at
index \(A\), has \(q=r+1-A-h\) remaining indices after it.
Its nonzero block entry \(\xi\) is multiplied by
\[
                  \tau^2\eta^{2q+h-1}.
\]
Except for a terminal singleton, its original density
factor is
\[
                 |\xi|^{\,2h(2q+k+h-2)}.
\]
For a terminal singleton the factor is instead
\((|\xi|^2-TT')^{k-1}\); it has the same scaling
because \(T,T'\) are both multiplied by \(\tau^2\).
The two elementary exponent identities are
\[
 \begin{split}
 4\sum_{v=0}^{h-1}(2(q+v)+k-1)
       &=4h(2q+k+h-2),\\
 4\sum_{v=0}^{h-1}(q+v)(2(q+v)+k-1)
       &=(2q+h-1)\,2h(2q+k+h-2)
                  +4\,\frac{h(h^2-1)}6 .
 \end{split}                                           \tag{CT12}
\]
They follow by summing \(1,v,v^2\).
The last term is four times the sum of \(j-i\) over
pairs in that block. Multiplying over all blocks gives
\[
       \mathrm{ptWeight}_{L_1}(\gamma_1)
                =W\eta^{-4S_\gamma}
                                \mathrm{ptWeight}_L(\gamma).
                                                        \tag{CT13}
\]
Since
\(\alpha_\gamma^2\eta^{-2d_\gamma}=\eta^{-4S_\gamma}\),
this proves \eqref{eq:character-weight}. This computation
uses the actual block density, and does not assume that
internal leading minors are nonzero.

On the big cell the root pattern is empty. Hence
(CT11) has scalar one. Its radial variables change by
\[
 (q_0,\ldots,q_{r-1},\zeta)
       \longmapsto
 (t_0^2q_0,\ldots,t_{r-1}^2q_{r-1},\tau^2\zeta).
                                                        \tag{CT14}
\]
This real change of variables has Jacobian
\[
       \left(\prod_{j<r}t_j^4\right)\tau^4
                    =\tau^{4s}\eta^{2rs}=J.
                                                        \tag{CT15}
\]
The original radial weight has factor \(W\):
its exponents at \(q_j\) are \(2(2(r-j)+k-1)\),
and its terminal factor has exponent \(k-1\).
This proves \eqref{eq:character-norm} by Tonelli and
the exact change of variables. The excluded zero
coordinates and the Gram-boundary circle have measure
zero, also when \(k=1\); if \(r=0\), the first product
is empty.

Finally every leading minor of order \(\ell\) in
(CT2) is multiplied by \(\prod_{j<\ell}t_j^2\).
In
\[
       \mathcal G=|\Delta_{r+1}|^2-TT'|\Delta_r|^2
\]
both terms are therefore multiplied by \(J\), proving
\eqref{eq:character-gram}. Multiplying
\eqref{eq:character-orb} and \eqref{eq:character-weight}
and comparing with \eqref{eq:character-norm} cancels
the entire positive factor \(W\). The remaining factor
is
\[
              J\eta^{2d_\gamma}=a^{2rs-2d_\gamma},
\]
which proves (CT4). All cancellations concern positive
finite scalars; zero orbital values and zero radial
norms require no separate division.
\end{proof}

\begin{corollary}[Uniform cost of the actual-character return]
\label{cor:archimedean-character-return}
Suppose \(b>1\), \(a,a^{-1}\le b^{C_h}\), and the retained
standard-character family satisfies
\[
 \mathrm{ptWeight}_{L_1}(\gamma_1)
      \mathrm{orbSq}_{L_1}(F_1,\gamma_1)
 \le b^B
       \max\{1,\mathcal G_{L_1}(\gamma_1)^{-1}\}^{p}
                           \mathcal N_{L_1}(F_1),
 \qquad p\ge0 .
\]
Then the original datum satisfies
\[
 \mathrm{ptWeight}_L(\gamma)\mathrm{orbSq}_L(F,\gamma)
 \le b^{B+2C_hrs(1+p)}
       \max\{1,\mathcal G_L(\gamma)^{-1}\}^{p}
                                      \mathcal N_L(F).
                                                        \tag{CT16}
\]
The parameter, compact types, and spectral reserve are
unchanged. In particular \(p=O(s^2)\) gives an additional
exponent \(O(C_hs^4)\), with no dependence on \(k\).
\end{corollary}
\begin{proof}
The root pattern consists of pairs of indices in the
same block, so
\[
             0\le2d_\gamma\le r(r+1)=rs.
\]
Thus \(0\le2rs-2d_\gamma\le2rs\). Also
\[
 \max\{1,(a^{2rs}\mathcal G)^{-1}\}
       \le\max\{1,a^{-2rs}\}\max\{1,\mathcal G^{-1}\}.
\]
Insert these two bounds in (CT4). The inequalities on
\(a,a^{-1}\) imply \(C_h\ge0\), and give (CT16).
For the original local admissibility gate,
\[
 T_1(1/a)^{2r}=T,\qquad T'_1(1/a)^{2r}=T',
 \qquad T_1=Ta^{2r}\ge1,\quad T'_1=T'a^{2r}\ge1.
\]
These are exactly its original standardized size data
and its two lower-bound hypotheses. Thus this transport
does not enlarge the required spectral family or consume
any additional shifted-weight reserve.
\end{proof}

This proposition transports only a standard-character
pointwise estimate proved on the stated retained family.
It does not assert such an estimate for a larger family
or for a weaker spectral-reserve hypothesis.
